% !TeX root = ../drafts/10-isa-programming.tex
\section{ISA programming}\label{sec:isa-programming}

How the zkDSL compiler (\path{crates/lean_compiler}; surface language in \path{crates/lean_compiler/zkDSL.md}) programs the seven-instruction ISA. Two themes run through every pattern: \emph{write-once memory is an assertion mechanism} (a second write of a cell is an equality check, an unwritten cell is prover-chosen), and \emph{indices live in the exponent} (a logical index $i$ is the element $\gen^{i}$, so incrementing, offsetting, and address formation are single field multiplications).

\subsection{Hints}\label{sec:prog-hints}

The prover's nondeterminism enters through write-once cells that are still unset when execution needs them. Some are explicit advice (a fresh allocation pointer $\gen^{\mathrm{base}}$, a witness stream, a computed-advice builtin); others are back-solved by an instruction whose remaining operands are known. In every case the resulting word is unconstrained except by the instructions and memory equalities that consume it, so a program must check all security-relevant advice.

\subsection{Division and inequality}\label{sec:prog-div-ne}

The zkDSL's single-slash division $q=a/b$ uses write-once back-solving: it emits \texttt{MUL\_NATIVE} with $q$ as the one unset input and $a$ as the already-written output. Witness generation fills $q=a\,b^{-1}$, while the ordinary multiplication table proves $q b=a$. It costs one VM instruction. Division by zero is undefined. Double slash \texttt{//} remains compile-time integer floor division for sizes and indices; it is not a field operation.

The assertion \texttt{assert a != b} computes $x=a+b$ with \texttt{XOR}, takes a hinted $w=x^{-1}$, forms $p=x\,w$ with \texttt{MUL\_NATIVE}, and writes $1$ to $p$ with \texttt{SET\_CONSTANT}. Memory being write-once, the two writes to $p$ agree only if $p=1$, and $x=0$ forces $p=0$ whatever the hint. Three instructions, and no branch: the condition of a \texttt{JUMP} would be the $\E$-valued $x$, whereas every branch a program actually needs compares $\gen$-powers.

\subsection{Functions}\label{sec:prog-functions}

A frame is fp-relative with layout $[\mathit{retpc},\mathit{retfp},\mathrm{args},\mathrm{rets},\mathrm{locals}]$. A call: the callee frame pointer is hinted into a fresh cell $\mathit{nfp}$; the arguments and the caller's $\fp$ are stored through $\mathit{nfp}$ with \texttt{DEREF}s (modes \texttt{cell} and \texttt{fp}); a \texttt{DEREF} in mode \texttt{pc} stores the return address $\gen^{2}\cdot\pc$ (the instruction after the transfer); one \texttt{JUMP} enters the callee. Return is a single \texttt{JUMP} to $\loc{1}$ with frame pointer $\loc{\gen}$. Cost: $n_{\mathrm{args}}+n_{\mathrm{rets}}+4$ instructions per call.

\subsection{Loops}\label{sec:prog-loops}

A counted loop keeps its counter in the exponent: it starts at $\gen^{\mathrm{lo}}$, advances by one multiplication $i\gets\gen\cdot i$, and exits on reaching $\gen^{\mathrm{hi}}$. The body is lowered as a tail-recursive helper function whose exit test \emph{is} the recursive call's \texttt{JUMP} condition (the \texttt{XOR} $i+\gen^{\mathrm{hi}}$ is nonzero exactly while the loop must continue), so an iteration costs one call plus two instructions, with no is-zero gadget. Loop state is threaded through captured arguments or a heap buffer.

\subsection{Conditionals}\label{sec:prog-conditionals}

\texttt{if}/\texttt{else} on a field equality is one \texttt{XOR} plus one conditional \texttt{JUMP}: the taken jump goes to whichever block the test should \emph{not} fall into ($=$ falls into \emph{then}, $\neq$ into \emph{else}), so no negation is ever computed. Intra-function jump targets are bytecode addresses $\gen^{\mathrm{entry}+i}$, backpatched at layout.

A taken \texttt{JUMP} reloads $\fp$ from a cell, and the ISA has no $\fp$-read; a branching function therefore materializes its own $\fp$ once: a \texttt{DEREF} in mode \texttt{fp} writes $\fp$ into a fresh one-cell heap buffer, a \texttt{DEREF} in mode \texttt{cell} copies it back into the frame (2 instructions, amortized; in \texttt{main}, $\fp=\gen^{0}=1$ is already a constant). Bindings made inside a branch are local to it; branches communicate through write-once cells: only one branch executes, so both may write the \emph{same} cell.

\subsection{Range checks}\label{sec:prog-range-checks}

The range check \emph{in the exponent} proves $x\in\{\gen^{0},\dots,\gen^{k-1}\}$, i.e.\ $\log_{\gen}x<k$, in 3 cycles (leanVM's \texttt{DEREF} trick, transported to $\gen$-powers):
\begin{enumerate}
\item \texttt{DEREF} through $x$: the dereferenced address $x\cdot\gen^{0}$ rides the memory interaction, which balances only for the $2^{h}$ addresses $\gen^{0},\dots,\gen^{2^{h}-1}$ (\S\ref{sec:memchan}): the bus itself proves $x=\gen^{e}$ with $e<2^{h}$;
\item \texttt{MUL} $x\cdot y$ into a write-once cell holding the constant $\gen^{k-1}$: the runner back-solves the complement $y=\gen^{k-1-e}$ (\S\ref{sec:prog-hints}), and the double write asserts $x\cdot y=\gen^{k-1}$;
\item \texttt{DEREF} through $y$: proves $y=\gen^{f}$, $f<2^{h}$.
\end{enumerate}
Then $e+f\equiv k-1\pmod{2^{64}-1}$ with $e,f<2^{h}\le2^{32}$, so $e+f<2^{33}<2^{64}-1$ and the congruence is an equality over the integers; a ``negative'' $k-1-e$ would wrap to $\approx2^{64}\gg2^{h}$, so $e\le k-1$, for every memory size the prover may announce, provided $k\le2^{16}$ (the minimum, \S\ref{sec:vm}). The two \texttt{DEREF} targets are deferred touches; the constant cell is one amortized \texttt{SET} per distinct bound per frame.

\subsection{Match statements}\label{sec:prog-match}

leanVM dispatches a switch with a single jump to $\pc=a+b\cdot x$: the case blocks, padded to a common length $b$, sit consecutively from base $a$, and its integer arithmetic forms the affine address in one instruction. That computation does not transport to this VM. Addresses here are $\gen$-powers: the additive offset $a+j$ is the cheap product $\gen^{a}\cdot x$, but the \emph{scaling} $b\cdot j$ is the power $x^{b}$ ($\log_2 b$ squarings per dispatch), with every block still padded to the longest. Instead, the aligned region collapses to fixed two-instruction \emph{trampoline slots} and the real blocks move out of line: matching on $x=\gen^{j}$, $j\in[0,n)$, the dispatch jumps to
\[
  d \;=\; \gen^{T}\cdot x^{2} \;=\; \gen^{T+2j},
\]
slot $j$ of a table at bytecode base $T$, where a slot is $\texttt{SET}\ c=\gen^{\mathrm{block}_j}$ then $\texttt{JUMP}\ c$. The case blocks sit anywhere, unaligned and of any length; the scaling cost is the single squaring $x^{2}$. Two jumps per dispatch, $\approx7$ cycles independent of $n$.

The slots are two instructions rather than one because a \texttt{JUMP} reads its destination from a \emph{cell}: a one-instruction slot would need its target cell pre-written, i.e.\ $n$ \texttt{SET}s executed before every dispatch. (Other layouts exist and trade exactly this, e.g.\ a memory-resident table of the $n$ block addresses, read with one \texttt{DEREF} and taken with a \emph{single} jump, paying the \texttt{SET}s as setup; worthwhile for many repeated small matches. The compiler currently emits only the trampoline.)

As in leanVM, the matched value must be known to lie in $[0,n)$: a hinted $x$ must be range-checked first (\S\ref{sec:prog-range-checks}); otherwise the dispatch lands at an attacker-chosen program counter.
